\documentclass[11pt]{article}

\usepackage[margin=1in]{geometry}
\usepackage{amsmath,amssymb,amsthm,mathtools}
\usepackage{enumitem}
\usepackage[hidelinks]{hyperref}
\usepackage{xcolor}
\usepackage{booktabs}

\newtheorem{theorem}{Theorem}[section]
\newtheorem{lemma}[theorem]{Lemma}
\newtheorem{proposition}[theorem]{Proposition}
\newtheorem{corollary}[theorem]{Corollary}
\theoremstyle{definition}
\newtheorem{definition}[theorem]{Definition}
\newtheorem{conjecture}[theorem]{Conjecture}
\newtheorem{leanblock}[theorem]{Lean formalization block}
\theoremstyle{remark}
\newtheorem{remark}[theorem]{Remark}
\newtheorem{warning}[theorem]{Warning}

\newcommand{\PP}{\mathbb P}
\newcommand{\EE}{\mathbb E}
\newcommand{\one}{\mathbf 1}
\newcommand{\cC}{\mathcal C}
\newcommand{\Sur}{\operatorname{Sur}}
\newcommand{\con}{\leftrightarrow}
\newcommand{\ncon}{\not\leftrightarrow}
\newcommand{\KN}{Kozma--Nitzan}
\newcommand{\lean}[1]{\texttt{#1}}

\title{All Numbered Kozma--Nitzan Conjectures and Questions:\\
       a Source-Set Proof and Lean Formalization Blueprint}
\author{Prepared from the supplied notes and formal development}
\date{2 September 2026}

\begin{document}
\maketitle

\begin{center}
\fbox{\parbox{0.92\textwidth}{
\textbf{Result.}
All five numbered conjectures in Kozma--Nitzan---Conjectures
1, 2, 3, 4, and 6---and all four numbered questions---Questions
5, 7, 8, and 9---follow from a guarded, source-set extension of the
conditioned slack hierarchy.  That extension is proved below from the
same finite-product and BHK inputs used by the supplied development.  In
fact we prove stronger statements.  In Conjecture~4 one may use the
particular relay minimizing the \emph{unconditional} cluster mean, which
answers Question~7.  A finite-dimensional separation argument gives the
common coefficients requested in Question~5.  The strong pre-FKG theorem
answers Question~8, and conditioning on the star of the observer answers
Question~9.  Finally, Conjecture~6 remains true without either of its two
endpoint assumptions.  The central device is to allow a finite set of
vertices, rather than one vertex, to be the source in the avoided
first-relay inequality.

The full source-set wrapper described here has not yet been submitted to a
Lean kernel.  The probabilistic input it needs is already present in the
repository, including set clusters, set markers, two-set conditional
association, and the CSH peeling theorem.  Section~\ref{sec:lean} separates
the new declarations into small formalization blocks.
}}
\end{center}

\begin{abstract}
Let independent bond percolation be given on a finite weighted graph.  For
a finite source set $S$, write $C_S=\bigcup_{s\in S}C_s$.  We prove a
source-set, avoided version of the first-relay inequality.  If $Y$ is an
avoided set, $T$ is a relay set, $\Phi$ is increasing, and the relays are
ordered by the conditional means
\[
 m_x^Y=\EE[\Phi(C_x)\mid x\ncon Y],
\]
then
\[
 \EE[\Phi(C_S);\ S\ncon Y,\ S\con T]
 \geq \sum_{x\in T}\PP(J_x^{S,Y})m_x^Y,
\]
where $J_x^{S,Y}$ says that $x$ is the first relay met by $C_S$.
The proof is the CSH surplus-transfer and relay-peeling proof with a
set-valued source.  We give all of the quantities and the two inequalities
used in the induction.

A monotone projection then shows that if $a\in A$ minimizes
$\EE F(C_a)$, it is the particular relay which works in Conjecture~4.
Taking $F(K)=\one\{b\in K\}$ proves the strong pre-FKG inequality and hence
Conjectures~2, 1, and 3.  Finally take the source set $S=\{v,w\}$.  The union
$C_v\cup C_w$, its connection events, and its avoidance events are
unchanged by the state of the internal edge $e=\{v,w\}$.  The avoided
source-set inequality therefore transfers verbatim to the graph in which
$e$ is forced open, proving Conjecture~6 with the old minimizer $a$.

The same fixed-minimizer theorem resolves the remaining numbered
questions.  Convex separation produces one probability vector on $A$
which works simultaneously for every target $b$ (Question~5).  Questions
7 and 8 are respectively the fixed-minimizer and minimum forms of the
strong pre-FKG inequality.  For Question~9, reveal all pairs incident to
the observer; their open endpoints form a random source set in the graph
with the observer's star deleted, to which the source-set theorem applies
configuration by configuration.
An elementary finite-edge exhaustion extends the needed inequality to
countable graphs; combined with Kozma--Nitzan's Theorem~6, this also gives
$\PP_{p_c}(|C(0)|=\infty)=0$ on $\mathbb Z^d$ for every $d\geq2$.
\end{abstract}

\tableofcontents

\section{Model and exact targets}\label{sec:targets}

Let $V$ be a finite vertex set and let every unordered pair $e$ be open,
independently, with probability $p_e\in[0,1]$.  A pair of weight zero is
simply an absent edge.  The open cluster of $x$ is $C_x$.  For a finite
source set $S\subseteq V$, put
\begin{equation}\label{eq:setcluster}
 C_S:=\bigcup_{s\in S}C_s.
\end{equation}
Thus $S\con B$ means $C_S\cap B\ne\varnothing$, and $S\ncon B$ means
$C_S\cap B=\varnothing$.  We use the same symbols when one of the sets is
a singleton.  All sets in this document are finite.

A cluster functional is an increasing map $F:2^V\to\mathbb R$.  The
paper's slightly more intrinsic definition of a monotone cluster property
is equivalent on the support of a finite percolation measure: condition
(2) in that definition makes the value common to all roots of one cluster,
and an isotone function on the poset of realizable clusters extends to
$2^V$ by
\[
 F(K)=\max\{f(L):L\text{ is a realizable cluster and }L\subseteq K\}.
\]
An irrelevant constant can be inserted when the displayed set is empty.

We record the five numbered conjectures in the form used below.  We always
assume $A\ne\varnothing$.

\begin{conjecture}[\KN\ 1]
For vertices $o,b$ and a relay set $A$,
\begin{equation}\label{eq:C1}
 \PP(o\con b)\geq \PP(o\con A)\min_{x\in A}\PP(x\con b).
\end{equation}
\end{conjecture}

\begin{conjecture}[\KN\ 2]
One has
\begin{equation}\label{eq:C2}
 \PP(o\con b)\geq
 \min_{x\in A}\PP(o\con A,\ x\con b).
\end{equation}
The useful stronger form has $\PP(o\con b,o\con A)$ on the left.
\end{conjecture}

\begin{conjecture}[\KN\ 3]
For every $\varepsilon>0$ there is a $\delta>0$, independent of the graph
and of $|A|$, such that
\[
 \PP(o\con A)>1-\delta,
 \qquad \PP(x\con b)>1-\delta\quad(x\in A)
\]
imply $\PP(o\con b)>1-\varepsilon$.
\end{conjecture}

\begin{conjecture}[\KN\ 4]
For every increasing cluster functional $F$,
\begin{equation}\label{eq:C4}
 \EE[F(C_o);o\con A]
 \geq \min_{x\in A}\EE[F(C_x);o\con A].
\end{equation}
\end{conjecture}

For $e=\{v,w\}$, write $G/e$ for the weighting obtained by replacing
$p_e$ by $1$.  This is equivalent to contraction for connectivity events.

\begin{conjecture}[\KN\ 6]
Suppose $a\in A$ minimizes $\PP_G(x\con b)$ over $x\in A$ and
\[
 \PP_G(x\con b)\geq
 \PP_G(x\con A)\PP_G(a\con b),\qquad x\in\{v,w\}.
\]
Then
\begin{equation}\label{eq:C6}
 \PP_{G/e}(v\con b)\geq
 \PP_{G/e}(v\con A)\PP_{G/e}(a\con b).
\end{equation}
\end{conjecture}

\begin{conjecture}[\KN\ Question 5]
Do there exist coefficients $c_x\geq0$, depending on $G,o,A$ but not on
$b$, with $\sum_{x\in A}c_x=1$, such that simultaneously for every
$b\in V$,
\begin{equation}\label{eq:Q5}
 \PP(o\con b)\geq
 \sum_{x\in A}c_x\PP(o\con A,x\con b)?
\end{equation}
\end{conjecture}

The remaining candidate questions all ask whether a specified relay $a$
satisfies
\begin{equation}\label{eq:Qcandidate}
 \PP(o\con b,o\con A)\geq\PP(o\con A,a\con b).
\end{equation}

\begin{conjecture}[\KN\ Question 7]
Does \eqref{eq:Qcandidate} hold when $a$ minimizes
$\PP(x\con b)$ over $x\in A$?
\end{conjecture}

\begin{conjecture}[\KN\ Question 8]
Does \eqref{eq:Qcandidate} hold when $a$ minimizes
$\PP(o\con A,x\con b)$ over $x\in A$?
\end{conjecture}

\begin{conjecture}[\KN\ Question 9]
Let $H$ be obtained from $G$ by deleting every pair incident to $o$.
Does \eqref{eq:Qcandidate}, with probabilities evaluated in $G$, hold
when $a$ minimizes $\PP_H(x\con b)$ over $x\in A$?
\end{conjecture}

There is no Conjecture~5 and no Question~6 in the paper.  Its complete
numbered inventory is
\[
 \text{Conjectures }1,2,3,4,6,
 \qquad \text{Questions }5,7,8,9.
\]
The results below prove every one of these statements.  The unnumbered
strong form (3) of Conjecture~2 is also proved.

\section{Avoided first-relay quantities}\label{sec:definitions}

Fix an avoided set $Y\subseteq V$.  For a vertex $x$ put
\[
 D_Y(x):=\{x\ncon Y\},\qquad d_Y(x):=\PP(D_Y(x)).
\]
For a source set $S$, define analogously
\[
 D_Y(S):=\{S\ncon Y\}=\bigcap_{s\in S}D_Y(s).
\]
Let $T\subseteq V\setminus Y$ be a relay set and let
$r:T\to\{1,\ldots,|T|\}$ be a bijective rank.  For an increasing
functional $\Phi$, assume first that $d_Y(x)>0$ for every $x\in T$ and put
\begin{equation}\label{eq:condmean}
 m_x^Y:=\frac{\EE[\Phi(C_x);D_Y(x)]}{d_Y(x)}.
\end{equation}
The rank is called compatible if
\begin{equation}\label{eq:compatible}
 r(x)<r(z)\quad\Longrightarrow\quad m_x^Y\leq m_z^Y.
\end{equation}

The first-relay pattern of $x$, seen from $S$ while $Y$ is avoided, is
\begin{equation}\label{eq:pattern}
 J_x^{S,Y}:=D_Y(S)\cap\{S\con x\}
 \cap\bigcap_{\substack{z\in T\\r(z)<r(x)}}\{S\ncon z\}.
\end{equation}
The patterns are disjoint and
\begin{equation}\label{eq:partition}
 \bigsqcup_{x\in T}J_x^{S,Y}
 =D_Y(S)\cap\{S\con T\}.
\end{equation}

Define the avoided source surplus
\begin{equation}\label{eq:surplus}
 \Sur_{S,Y}(T;r,\Phi):=
 \EE[\Phi(C_S);D_Y(S),S\con T]
 -\sum_{x\in T}\PP(J_x^{S,Y})m_x^Y.
\end{equation}
For a singleton source we abbreviate $\Sur_{x,Y}:=\Sur_{\{x\},Y}$.
Notice that no connectedness of $C_S$ is asserted or needed.

\section{The source-set lift of the conditioned slack hierarchy}
\label{sec:sourceset}


The extension needed below is not obtained by wiring the source vertices.
Wiring would turn $C_S$ into one connected cluster, but would also change
the law of a relay cluster whenever that relay reaches $S$.  There is a
second, subtler trap.  From
\[
 z\ncon B,\qquad z\con S
\]
one cannot conclude $S\ncon B$: a different component generated by $S$
may hit $B$.  Thus an unguarded set-observer version of CSH is not enough.
The correct extension retains the full source-avoidance guard at every
evaluation.

\subsection{Guarded source forms}

For finite vertex sets $R,X,Y$, write
\begin{equation}\label{eq:guard-event}
 \mathsf H_{X\mid Y}(R):=\{R\ncon Y,\ R\con X\}.
\end{equation}
The elementary identity driving the construction is the disjoint union
\begin{equation}\label{eq:guard-split}
 \mathsf H_{X\cup\{d\}\mid Y}(R)
 =\mathsf H_{X\mid Y}(R)\ \sqcup\
   \mathsf H_{\{d\}\mid Y\cup X}(R).
\end{equation}
Indeed, if $C_R$ has not met $X$, then its first new contact with
$X\cup\{d\}$ is a contact with $d$, and in that case $C_R$ avoids
$Y\cup X$.  Notice that \eqref{eq:guard-split} keeps track of every
component in $C_R$.

For a vertex owner $x$, an avoided set $Y$, a finite source $R$, and an
increasing functional $\psi$, put
\begin{align}
 d_x(Y)&:=\PP(x\ncon Y),\label{eq:guard-denom}\\
 \Gamma_{x,Y}^{\psi}(R)
 &:=d_x(Y)\EE[\psi(C_x);\mathsf H_{\{x\}\mid Y}(R)]\notag\\
 &\hspace{1.5cm}-\EE[\psi(C_x);x\ncon Y]
       \PP(\mathsf H_{\{x\}\mid Y}(R)).\label{eq:guardGamma}
\end{align}
When $R=\{u\}$, the guard and the contact imply that $C_u=C_x$ and
$x\ncon Y$; hence \eqref{eq:guardGamma} is exactly the ordinary
denominator-free CSH covariance evaluated at $u$.

Let ${\bf d}=(d_1,\ldots,d_\ell)$ be a list of distinct decoys and set
\[
 B_j=Y\cup\{x,d_i:i<j\},\qquad
 B_*=Y\cup\{x,d_1,\ldots,d_\ell\}.
\]
The guarded decoy profile and observer price are
\begin{align}
 c_j(R)&=
 \frac{\PP(\mathsf H_{\{d_j\}\mid B_j}(R))}
      {\PP(d_j\ncon B_j)},\label{eq:guard-decoy}\\
 p_{O,v}&=
 \frac{\PP(\mathsf H_{\{v\}\mid B_*}(O))}
      {\PP(v\ncon B_*)}.
 \label{eq:guard-price}
\end{align}
For a function $h$ on finite source sets define
\begin{equation}\label{eq:guard-slack}
 L_0 h=h,\qquad
 (L_j h)(R)=(L_{j-1}h)(R)-c_j(R)(L_{j-1}h)(\{d_j\}).
\end{equation}
All evaluations other than the primary one at $O$ are singletons.  In
particular, for singleton $O=\{o\}$ all guards in
\eqref{eq:guard-decoy}--\eqref{eq:guard-price} collapse by connectivity,
and these are precisely the constants of the CSH already in the project.

\begin{theorem}[guarded-source conditioned slack hierarchy]
\label{thm:guardCSH}
Assume first that all edge weights lie in $(0,1)$ and that the sets
$Y$, $O$, $\{v\}$, $\{x\}$, and $\{d_1,\ldots,d_\ell\}$ are pairwise
disjoint (the $d_j$ are distinct).  For every increasing $\psi$,
\begin{equation}\label{eq:guardCSH}
 \boxed{\quad
 L_\ell\Gamma_{x,Y}^{\psi}(O)
 -p_{O,v}L_\ell\Gamma_{x,Y}^{\psi}(\{v\})\geq0.
 \quad}
\end{equation}
\end{theorem}

\begin{proof}
We give every additional identity needed beyond the existing CSH proof.
Throughout this proof a covariance is unnormalised: thus
\(\operatorname{cov}_D(f,g)=\PP(D)\EE[fg;D]-\EE[f;D]\EE[g;D]\).
Every $\Gamma$ in the asserted margin annihilates constants.  Since the
state space is finite, we may therefore add a constant to $\psi$ and
assume throughout the proof that $\psi\geq0$.

\emph{1. The source kernel and guarded exchange.}
Work in an induced finite world $U$.  For a possible cluster $K$ of a
vertex $v$, define
\begin{equation}\label{eq:alpha-kernel}
 \alpha_{R,B}^{U}(K)=
 \one\{K\cap R\ne\varnothing\}\,
 \PP_{U\setminus K}(C_{R\setminus K}\cap B=\varnothing).
\end{equation}
Boundary pairs meeting $K$ are deleted in $U\setminus K$.  Directly from
this definition,
\begin{equation}\label{eq:alpha-properties}
 K\subseteq K'\Longrightarrow
 \alpha_{R,B}^{U}(K)\leq\alpha_{R,B}^{U}(K'),
 \qquad B\subseteq B'\Longrightarrow
 \alpha_{R,B'}^{U}(K)\leq\alpha_{R,B}^{U}(K).
\end{equation}
For the first implication, the indicator can only turn on, the residual
source shrinks, and the residual graph loses possible paths to $B$.
Exploration of $C_v$ gives the exact tower identity
\begin{equation}\label{eq:alpha-tower}
 \EE_U[\alpha_{R,B}^{U}(C_v);v\ncon B]
 =\PP_U(R\ncon B,R\con v).
\end{equation}
In particular, exploration of $C_x$ gives
\begin{equation}\label{eq:gamma-alpha}
 \Gamma_{x,Y}^{\psi}(R)
 =d_x(Y)\EE[\psi(C_x)\alpha_{R,Y}(C_x);x\ncon Y]
  -\EE[\psi(C_x);x\ncon Y]
   \EE[\alpha_{R,Y}(C_x);x\ncon Y].
\end{equation}
Thus a guarded contact can be averaged to an increasing function of one
cluster.  We use this identity below for the exchange estimate.  It is
important, however, not to feed $\alpha_{R,Y}^{V}$ unchanged into the
old one-cluster Gibbs reduction: after $C_Y$ is exposed, that fixed kernel
is not the residual-world contact indicator.  The correct reduction is
the two-cluster calculation in part~4 below.

Fix a marker set $X$ with $x\in X$ and $v\notin X$, and put
\[
 E_{X,R}(N)=\PP_U(R\ncon X\cup N,R\con v),\qquad
 M_X(N)=\PP_U(v\ncon X\cup N).
\]
If $N\cap N'=\varnothing$ and $N,N'$ are disjoint from $X$, then
\begin{equation}\label{eq:guard-exchange}
 E_{X,R}(N)M_X(N')
 \leq M_X(N\cup N')E_{X,R}(\varnothing).
\end{equation}
Here is a complete derivation.  The one-cluster BHK core says, for
increasing nonnegative $f,g$,
\[
 \EE[f(C_v);v\ncon A]\EE[g(C_v);v\ncon A']
 \leq \EE[f(C_v)g(C_v);v\ncon A\cap A']
       \PP(v\ncon A\cup A').
\]
Apply it with avoided sets $X\cup N$ and $X\cup N'$, increasing function
$K\mapsto\alpha_{R,X\cup N}^{U}(K)$, and the constant function $1$.
Because the two additional sets are disjoint, their intersection is
$X$ and their union is $X\cup N\cup N'$.  Thus
\[
 \EE[\alpha_{R,X\cup N}(C_v);v\ncon X\cup N]M_X(N')
 \leq
 \EE[\alpha_{R,X\cup N}(C_v);v\ncon X]M_X(N\cup N').
\]
The first expectation is $E_{X,R}(N)$ by
\eqref{eq:alpha-tower}.  By \eqref{eq:alpha-properties}, the second is at
most the same expectation with $X\cup N$ replaced by $X$, namely
$E_{X,R}(\varnothing)$.  This proves \eqref{eq:guard-exchange}.  Cases in
which the displayed disjointness fails reduce either to the same statement
after deleting redundant vertices or to a zero factor.

\emph{2. The guarded two-source inequality.}
For a nonnegative world functional $h$ write, with
$U_N=U\setminus C_N$,
\begin{align*}
 Y_h^x(N)&=\EE_U[\one\{x\ncon N\}h(U_N)],\\
 X_{R,h}^x(N)&=\EE_U[\one\{x\ncon N,R\ncon N\}
                 q_R(U_N)h(U_N)],\\
 q_R(W)&=\frac{\PP_W(R\ncon X,R\con v)}
                    {\PP_W(v\ncon X)}.
\end{align*}
The one-source stopped-Harris estimate already proved in the project is
\begin{equation}\label{eq:guard-star}
 Y_h^x(N)M_X(\varnothing)\leq M_X(N)h(U),
 \qquad N\longmapsto Y_h^x(N)\text{ is antitone},
\end{equation}
for
$h(W)=\operatorname{Cov}_W(\psi(C_x),\one\{v\con X\})$.
The guarded analogue of \lean{CovTau.metaA2\_of\_star} is
\begin{equation}\label{eq:guard-A2}
 E_{X,R}(N)Y_h^x(N')
 \leq M_X(N\cup N')X_{R,h}^x(N\cap N').
\end{equation}
For $N\cap N'=\varnothing$, multiply \eqref{eq:guard-star} by
$E_{X,R}(N)$ and use \eqref{eq:guard-exchange}.  Since
$E_{X,R}(\varnothing)=q_R(U)M_X(\varnothing)$, cancellation of the
positive factor $M_X(\varnothing)$ gives
\[
 E_{X,R}(N)Y_h^x(N')
 \leq M_X(N\cup N')q_R(U)h(U),
\]
which is \eqref{eq:guard-A2} at the empty intersection.  If that factor
is zero, both $E_{X,R}(N)$ and the left side vanish, so the same conclusion
holds without division.

For completeness we give the nonempty-intersection step rather than
referring to the analogous step in the ordinary META--A2 proof.  Put
$Z=N\cap N'\ne\varnothing$ and $U'=U\setminus Z$.  The cases
$R\cap Z\ne\varnothing$, $x\in Z$, or $v\in Z$ have a zero factor and
are immediate.  Otherwise expose the set $s\subseteq U'$ of vertices
having at least one open pair to $Z$.  The indicators defining $s$ use
disjoint blocks of pair coordinates, so its law is a product law.  With
$N_0=N\setminus Z$ and $N'_0=N'\setminus Z$, source-cluster exploration
gives the following four exact residual functions:
\[
 \begin{array}{ll}
 e(s)=E_{X,R}^{U'}(N_0\cup s),
 &y(s)=Y_{h}^{x,U'}(N'_0\cup s),\\[2mm]
 m(s)=M_X^{U'}(N_0\cup N'_0\cup s),
 &\chi(s)=X_{R,h}^{x,U'}(s).
 \end{array}
\]
The extra source guard has an exact recursion here: on the nonzero
branch, $R\ncon N$ in $U$ is equivalent to
$R\ncon(N_0\cup s)$ in $U'$.  Thus no residual avoidance event has been
factored.

For two possible neighbour sets $s,t\subseteq U'$, define the
cross-deleted arguments
\[
 P=s\cup(N_0\setminus t),\qquad
 P'=t\cup(N'_0\setminus s).
\]
Both $E_{X,R}$ and $Y_h^x$ are antitone.  Since
$P\subseteq N_0\cup s$ and $P'\subseteq N'_0\cup t$, nonnegativity and
the induction hypothesis in $U'$ give
\[
 e(s)y(t)
 \leq E_{X,R}^{U'}(P)Y_h^{x,U'}(P')
 \leq M_X^{U'}(P\cup P')X_{R,h}^{x,U'}(P\cap P').
\]
As $N_0\cap N'_0=\varnothing$, elementary set algebra gives
\[
 P\cup P'=N_0\cup N'_0\cup(s\cup t),
 \qquad P\cap P'=s\cap t.
\]
Consequently
\[
 e(s)y(t)\leq m(s\cup t)\chi(s\cap t).
\]
All four functions are nonnegative, and the product density of the
neighbour set satisfies
$\pi(s)\pi(t)=\pi(s\cap t)\pi(s\cup t)$.  The
Ahlswede--Daykin four-functions theorem therefore integrates this
pointwise inequality to \eqref{eq:guard-A2}.  Together with the empty
intersection case, this proves \eqref{eq:guard-A2} by strong induction
on $|U|$.

Take $N=N'=Y$ in \eqref{eq:guard-A2}.  With
$U_Y=V\setminus C_Y$ it says
\begin{multline}\label{eq:guard-Htw}
 \PP(R\ncon X\cup Y,R\con v)
 \EE[\one\{x\ncon Y\}h(U_Y)]\\
 \leq \PP(v\ncon X\cup Y)
 \EE[\one\{x\ncon Y,R\ncon Y\}q_R(U_Y)h(U_Y)].
\end{multline}

We also need the following guarded version of the four-point transfer.
Inside any fixed world put
\[
 h_R(U)=\operatorname{Cov}_U(\psi(C_x),\one\{R\con X\}),
 \qquad h_v(U)=\operatorname{Cov}_U(\psi(C_x),\one\{v\con X\}).
\]
Then
\begin{equation}\label{eq:guard-fourpoint}
 h_R(U)\geq q_R(U)h_v(U).
\end{equation}
Here is a proof, included because the superficially stronger coefficient
$\PP(R\con v\mid v\ncon X)$ is in general false for a disconnected
source $R$.  Put
\[
 A=\{R\con X\},\quad B=\{R\con v\},\quad
 D=\{v\ncon X\},\quad J=\{R\ncon X,R\con v\}.
\]
The increasing event $A\cup B$ and Harris give, with
$m=\EE_U\psi(C_x)$,
\begin{equation}\label{eq:guard-fourpoint-harris}
 h_R(U)\geq m\PP_U(J)-\EE_U[\psi(C_x);J],
\end{equation}
because $(A\cup B)\setminus A=J$.  On $D$ the event $J$ is
\[
 \{C_v\cap R\ne\varnothing\}\cap\{C_X\cap R=\varnothing\}.
\]
It is increasing in $C_v$ and decreasing in $C_X$, whereas
$\psi(C_x)$ is constant in the first cluster and increasing in the
second.  The two-set BHK conditional association theorem therefore gives
\begin{equation}\label{eq:guard-fourpoint-bhk}
 \PP_U(D)\EE_U[\psi(C_x);J]
 \leq \PP_U(J)\EE_U[\psi(C_x);D].
\end{equation}
Since
$h_v(U)=\PP_U(D)m-\EE_U[\psi(C_x);D]$, equations
\eqref{eq:guard-fourpoint-harris}--\eqref{eq:guard-fourpoint-bhk}
prove \eqref{eq:guard-fourpoint}; the zero-denominator case is again
immediate.

Divide
\eqref{eq:guard-Htw} by its positive denominator and combine it with
\eqref{eq:guard-fourpoint}.  The resulting horizontal margin is
\begin{equation}\label{eq:guard-horizontal}
 \EE[\one\{x\ncon Y,R\ncon Y\}h_R(U_Y)]
 -p_{R,v}\EE[\one\{x\ncon Y\}h_v(U_Y)]\geq0.
\end{equation}

\emph{3. The repaired guarded unfolding.}
We spell out the point at which the guard changes the existing
\lean{CSH.slForm\_jn} calculation.  Fix an outer configuration, let $Z$
be the already exposed cluster of $Y$, and work in the residual graph.
For its connectivity relation write
\[
 j_{S\mid Z}(R)=\one\{R\cap Z=\varnothing,\ R\con S\},\qquad
 a_{S,d\mid Z}(R)=\one\{R\cap Z=\varnothing,\ R\ncon S,\ R\con d\}.
\]
Then \eqref{eq:guard-split} is the pointwise identity
\begin{equation}\label{eq:guard-j-split}
 j_{S\cup\{d\}\mid Z}(R)=j_{S\mid Z}(R)+a_{S,d\mid Z}(R).
\end{equation}
If $e_{d,Z}=\one\{d\notin Z\}$, the singleton dichotomy is
\begin{equation}\label{eq:guard-singleton-dichotomy}
 j_{S\mid Z}(\{d\})=e_{d,Z}-a_{S,d\mid Z}(\{d\}).
\end{equation}
Consequently one level step, with any coefficient function $c$, is
\begin{multline}\label{eq:guard-level-step}
 j_{S\mid Z}(R)-c(R)j_{S\mid Z}(\{d\})\\
 =j_{S\cup\{d\}\mid Z}(R)
  -\bigl(a_{S,d\mid Z}(R)-c(R)a_{S,d\mid Z}(\{d\})\bigr)
  -c(R)e_{d,Z}.
\end{multline}
The last term is constant in the residual configuration and therefore
has zero covariance.  Iterating \eqref{eq:guard-level-step} is the exact
guarded replacement for \lean{CSH.slForm\_jn}; it produces the marker
$S=\{x,d_1,\ldots,d_\ell\}$ and one centred bracket for each decoy.

We now record the extra exploration which is essential for a set-valued
source.  It is useful to distinguish edge and vertex clusters.  Write
$\cC_R(\omega)$ for the union of the open \emph{edge} clusters rooted in
$R$, and put
\[
 V_R(F):=R\cup\{u:\exists e\in F,\ u\in e\},
 \qquad C_R(\omega)=V_R(\cC_R(\omega)).
\]
Thus isolated source vertices are retained.  In the Lean notation let
$g$ be the edge-cluster reading of $\psi$ and set
\[
 \Theta(\omega):=\lean{CSH.resid}(w,x,Y,g,\omega),\qquad
 \widehat\Phi(K):=\lean{CSH.phiFun}(w,x,Y,g,K).
\]
The correctly typed functional on an edge cluster of $d$ is
\begin{equation}\label{eq:typed-phi}
 \Psi_d(F):=\widehat\Phi\bigl(\{d\}\cup
                    \{u:\exists e\in F,\ u\in e\}\bigr),
 \qquad
 \Psi_d(\cC_d)=\widehat\Phi(C_d).
\end{equation}
The existing lemmas \lean{CSH.phiFun\_nonneg} and
\lean{CSH.phiFun\_mono} say that $\widehat\Phi$, and hence $\Psi_d$, is
nonnegative and increasing.

Let $S$ contain $x$, let $d\notin S\cup Y$, and put $B=S\cup Y$.  For an
arbitrary finite source $R$ abbreviate
\[
 E_R:=\mathsf H_{\{d\}\mid B}(R)
     =\{C_R\cap B=\varnothing,\ d\in C_R\},
 \qquad E_d:=\{d\ncon B\}.
\]
The correct set-cluster Markov identity is
\begin{equation}\label{eq:guard-markov-a}
 \EE[\Theta;E_R]=-\EE[\widehat\Phi(C_R);E_R].
\end{equation}
Here is a finite-sum proof, included to make the replacement formalizable.
For an edge set $F$ let
\[
 \iota(F)=\one\{V_R(F)\cap B=\varnothing,\ d\in V_R(F)\},
\]
and let $\overline F_R$ be the set of pairs meeting $V_R(F)$.  Part~(a)
of Lemma $\Phi$, in residual form, gives
\[
 \sum_\eta {\rm wt}(\eta)\,
   \Theta(\eta\setminus\overline F_R)
 =-\widehat\Phi(V_R(F)).
\]
Apply the set-cluster conditioning identity
\lean{CSH.set\_sum\_cond\_sdiff\_off} to the kernel
$(F,\beta)\mapsto\iota(F)\Theta(\beta)$.  On
$\{\iota(\cC_R)=1\}$ the whole union $C_R$ avoids
$\{x\}\cup Y\subseteq B$, so deleting all pairs meeting $C_R$ changes
neither $C_x$, $C_Y$, nor the residual $\Theta$.  The left side of the
conditioning identity is therefore the left side of
\eqref{eq:guard-markov-a}; the displayed residual identity turns its
right side into the right side of \eqref{eq:guard-markov-a}.  The point
is that the exploration is of the \emph{whole} source union $C_R$, not
only the one component $C_d$.

On $E_R$ one has $C_d\subseteq C_R$.  Define the unnormalised correction
\begin{equation}\label{eq:guard-error}
 \operatorname{Err}_{d,B}(R):=
 \EE\!\left[
   \widehat\Phi(C_R)-\Psi_d(\cC_d);E_R\right].
\end{equation}
Monotonicity of $\widehat\Phi$ gives
\begin{equation}\label{eq:guard-error-sign}
 \operatorname{Err}_{d,B}(R)\geq0.
\end{equation}
For every vertex $u$,
\begin{equation}\label{eq:guard-error-singleton}
 \operatorname{Err}_{d,B}(\{u\})=0,
\end{equation}
because on $E_{\{u\}}$ the connection $u\con d$ gives $C_u=C_d$.
Impossible events and the case $u=d$ are included.

Equations \eqref{eq:guard-markov-a}--\eqref{eq:guard-error} give
\[
 \EE[\Theta;E_R]
 =-\EE[\Psi_d(\cC_d);E_R]-\operatorname{Err}_{d,B}(R),
\]
whereas the singleton source gives the exact identity
\begin{equation}\label{eq:guard-markov-d}
 \EE[\Theta;E_d]
 =-\EE[\Psi_d(\cC_d);E_d].
\end{equation}
Put $e_d=\PP(E_d)$ and $q_R=\PP(E_R)$.  The division-free form, valid
also when $e_d=0$, is
\begin{equation}\label{eq:guard-one-decoy-cross}
 e_d\EE[\Theta;E_R]-q_R\EE[\Theta;E_d]
 =-\Gamma_{d,B}^{\widehat\Phi}(R)
  -e_d\operatorname{Err}_{d,B}(R).
\end{equation}
Under the present nondegeneracy assumptions $e_d>0$.  With
$c(R)=q_R/e_d$, and after the $C_Y$ Markov merge, the two $a$-terms are
the indicators of $E_R$ and $E_d$.  Hence the corrected one-decoy
identity is
\begin{equation}\label{eq:guard-one-decoy}
 \EE\!\left[\Theta
   \bigl(a_{S,d\mid C_Y}(R)-c(R)a_{S,d\mid C_Y}(\{d\})\bigr)\right]
 =-e_d^{-1}\Gamma_{d,B}^{\widehat\Phi}(R)
  -\operatorname{Err}_{d,B}(R).
\end{equation}

We now make every quantity in the unfolding explicit.  Let $\mu$ be the
original product law, put
\[
 D=\{x\ncon Y\},\qquad \delta=\PP(D)>0,
\]
and, for an outer configuration $\omega\in D$, let
$Z_\omega=C_Y(\omega)$.  Denote by $\mu_\omega$ the product law obtained
from $\mu$ by setting to zero every pair meeting $Z_\omega$.  Thus
$\mu_\omega$ is the residual-world law after the $Y$-cluster has been
exposed.  For a residual configuration $\eta$ define
\begin{align*}
 j^\omega_S(R;\eta)
   &=\one\{R\cap Z_\omega=\varnothing,\ C_R(\eta)\cap S\ne\varnothing\},\\
 a^\omega_{S,d}(R;\eta)
   &=\one\{R\cap Z_\omega=\varnothing,\ C_R(\eta)\cap S=\varnothing,
                    \ d\in C_R(\eta)\},\\
 \kappa^\omega_S(R)
   &=\operatorname{Cov}_{\mu_\omega}
       \bigl(\psi(C_x),j^\omega_S(R;\mathord\cdot)\bigr).
\end{align*}
The cluster symbols in these three formulas are evaluated in the residual
configuration.  These are precisely $j_{S\mid Z_\omega}$,
$a_{S,d\mid Z_\omega}$, and their world covariance from
\eqref{eq:guard-level-step}.

We shall also need tail level forms.  For $0\leq j\leq k\leq\ell$ put
\[
 L_{j,j}h=h,\qquad
 (L_{j,k}h)(R)=(L_{j,k-1}h)(R)
       -c_k(R)(L_{j,k-1}h)(\{d_k\}),
 \qquad L_{>j}:=L_{j,\ell}.
\]
Thus $L_{>0}=L_\ell$, while $L_{>j}$ uses exactly the later decoys
$d_{j+1},\ldots,d_\ell$.  Write
\[
 X=\{x,d_1,\ldots,d_\ell\},\qquad p=p_{O,v}.
\]
The unnormalised averaged within-world and horizontal margins are
\begin{align}
 \mathcal W_{x,Y,{\bf d}}^{O,v}(\psi)
   &:={\int_D}\Bigl[
       (L_{>0}\kappa^\omega_{\{x\}})(O)
       -p(L_{>0}\kappa^\omega_{\{x\}})(\{v\})\Bigr]d\mu(\omega),
       \label{eq:guard-W-def}\\
 \mathcal H_{x,Y,{\bf d}}^{O,v}(\psi)
   &:={\int_D}\Bigl[
       \kappa^\omega_X(O)-p\kappa^\omega_X(\{v\})\Bigr]d\mu(\omega).
       \label{eq:guard-H-def}
\end{align}
By cluster exploration, the second quantity is exactly
\[
 \EE[\one\{x\ncon Y,O\ncon Y\}h_O(U_Y)]
 -p\EE[\one\{x\ncon Y\}h_v(U_Y)],
\]
the left side of \eqref{eq:guard-horizontal} with $R=O$.  If
$v\in C_Y$, its residual-world contact indicator is zero.  Hence
$\mathcal H\geq0$.

For the $j$-th decoy set
\[
 S_j=\{x,d_i:i<j\},\qquad B_j=S_j\cup Y,\qquad
 e_j=\PP(d_j\ncon B_j),\qquad
 \operatorname{Err}_j=\operatorname{Err}_{d_j,B_j},
 \qquad \Psi_j=\Psi_{d_j}.
\]
The exact bridge between the residual-world calculation and the global
one-decoy identity is the following function of the source $R$:
\begin{align}
 I_j(R)
 &:={\int_D}\operatorname{Cov}_{\mu_\omega}\!\left(
       \psi(C_x),
       a^\omega_{S_j,d_j}(R;\mathord\cdot)
       -c_j(R)a^\omega_{S_j,d_j}(\{d_j\};\mathord\cdot)
       \right)d\mu(\omega)                                      \notag\\
 &=\EE[\Theta;\mathsf H_{\{d_j\}\mid B_j}(R)]
   -c_j(R)\EE[\Theta;d_j\ncon B_j]                              \notag\\
 &=-e_j^{-1}\Gamma_{d_j,B_j}^{\widehat\Phi}(R)
   -\operatorname{Err}_j(R).
 \label{eq:guard-world-global-bridge}
\end{align}
The second line is a kernel-valued $C_Y$ set-cluster Markov merge; unlike
a fixed-test conditioning statement, it retains the full guard
$R\ncon Y$.  For completeness, here is the finite-sum identity being
used.  Write $\operatorname{cut}(Z)$ for all pairs meeting $Z$ and put
\[
 \bar g(Z)=\sum_\eta {\rm wt}(\eta)\,
   \psi\bigl(C_x(\eta\setminus\operatorname{cut}(Z))\bigr).
\]
For every bounded two-argument test $F$ set-cluster conditioning gives
\begin{multline}\label{eq:guard-kernel-merge}
 \sum_\omega {\rm wt}(\omega)\one\{x\notin Z_\omega\}
 \sum_\eta {\rm wt}(\eta)
 \left[\psi\bigl(C_x(\eta\setminus\operatorname{cut}(Z_\omega))\bigr)
       -\bar g(Z_\omega)\right]
 F(Z_\omega,\eta\setminus\operatorname{cut}(Z_\omega))\\
 =\sum_\zeta {\rm wt}(\zeta)\Theta(\zeta)
 F(Z_\zeta,\zeta\setminus\operatorname{cut}(Z_\zeta)).
\end{multline}
This is the set-cluster tower identity applied to the kernel
$(Z,\beta)\mapsto\one\{x\notin Z\}
(\psi(C_x(\beta))-\bar g(Z))F(Z,\beta)$.  On the guarded branch,
$a^{\zeta}_{S_j,d_j}(R;\zeta\setminus\operatorname{cut}(Z_\zeta))$
is exactly the indicator of
$\mathsf H_{\{d_j\}\mid S_j\cup Y}(R)$ in $\zeta$.  This proves the
second line of \eqref{eq:guard-world-global-bridge}; its third line is
\eqref{eq:guard-one-decoy}.

Define the lower-level margin
\begin{equation}
 \mathcal M_j
 :=(L_{>j}\Gamma_{d_j,B_j}^{\widehat\Phi})(O)
   -p(L_{>j}\Gamma_{d_j,B_j}^{\widehat\Phi})(\{v\}).
 \label{eq:guard-M-def}
\end{equation}
There is no change of constants hidden here: for every $k>j$ the tail
profile in the guarded hierarchy with owner $d_j$, avoided set $B_j$,
and tail $(d_{j+1},\ldots,d_\ell)$ is $c_k$, and
\[
 B_j\cup\{d_j,d_{j+1},\ldots,d_\ell\}=B_*.
\]
Consequently its observer price is the same $p$, and $\mathcal M_j$ is
exactly the left side of \eqref{eq:guardCSH} for that lower-level datum.

Iterating the pointwise identity \eqref{eq:guard-level-step}, taking
world covariances, and then performing the final $O-pv$ subtraction gives
\begin{equation}
 \mathcal W
 =\mathcal H-
   \sum_{j=1}^{\ell}\Bigl[
      (L_{>j}I_j)(O)-p(L_{>j}I_j)(\{v\})\Bigr].
 \label{eq:guard-unfold-before-markov}
\end{equation}
All omitted terms here are constant in the residual configuration and
therefore have zero world covariance.  Moreover
\eqref{eq:guard-error-singleton} implies, by induction through the tail
steps, that $L_{>j}\operatorname{Err}_j=\operatorname{Err}_j$.
Substitution of \eqref{eq:guard-world-global-bridge} into
\eqref{eq:guard-unfold-before-markov} therefore gives the exact repaired
unfolding
\begin{equation}\label{eq:guard-unfold}
 \mathcal W_{x,Y,{\bf d}}^{O,v}(\psi)
 =\mathcal H_{x,Y,{\bf d}}^{O,v}(\psi)
  +\sum_{j=1}^{\ell}e_j^{-1}\mathcal M_j
  +\sum_{j=1}^{\ell}\operatorname{Err}_j(O).
\end{equation}
In particular,
\begin{equation}\label{eq:guard-unfold-lower}
 \mathcal W_{x,Y,{\bf d}}^{O,v}(\psi)
 \geq\mathcal H_{x,Y,{\bf d}}^{O,v}(\psi)
  +\sum_{j=1}^{\ell}e_j^{-1}\mathcal M_j.
\end{equation}

We prove Theorem~\ref{thm:guardCSH} by strong induction on $\ell$.  If
$\ell=0$, then $\mathcal W=\mathcal H\geq0$, and part~4 below turns this
within-world inequality into the level-zero theorem.  For $\ell>0$,
assume the theorem for every shorter decoy list.  Each $\mathcal M_j$
has only $\ell-j$ decoys and is nonnegative by that induction hypothesis,
applied to the increasing vertex functional $\widehat\Phi$; its correctly
typed edge-cluster reading in Lean is $\Psi_j$.  Since $e_j>0$ and every
error is nonnegative, \eqref{eq:guard-unfold} gives $\mathcal W\geq0$.
Part~4 turns this into the level-$\ell$ guarded CSH margin and closes the
induction.

\begin{remark}[regression test for the repaired identity]
The error term cannot be omitted, even with all weights strictly between
zero and one.  Take
$V=\{x,d,r_1,r_2\}$, give the pairs $\{x,r_2\}$ and $\{d,r_1\}$ weight
$1/2$ and every other pair weight $1/4$, and let
$Y=\varnothing$, $S=\{x\}$, $R=\{r_1,r_2\}$, and
$g(C_x)=\one\{r_2\in C_x\}$.  Exact enumeration gives
\[
 \EE[\Theta a_{S,d}(R)]=-\frac{1665}{16384},\qquad
 -\EE[\widehat\Phi(C_d);E_R]=-\frac{1017}{16384},
\]
and their discrepancy is
$-81/2048=-\operatorname{Err}_{d,B}(R)$.  Thus the former
$C_d$-only equality is false, while
\eqref{eq:guard-markov-a}--\eqref{eq:guard-one-decoy} give the exact
replacement.
\end{remark}

\emph{4. The two-cluster reduction, including the missing half-step.}
This is the point at which the source guard must be retained.  Recall
$D=\{x\ncon Y\}$ and $\delta=\PP(D)>0$.  Under the conditional law
$\PP(\,\cdot\mid D)$ put $A=C_x$ and $B=C_Y$, viewed as their vertex
sets, and define
\[
 b_R(A,B)=\one\{A\cap R\ne\varnothing\}\one\{B\cap R=\varnothing\}.
\]
Then $b_R(C_x,C_Y)$ is exactly the indicator of
$\mathsf H_{\{x\}\mid Y}(R)$.  Apply the level form and the final observer
subtraction to $R\mapsto b_R(A,B)$, and call the resulting bounded test
$\mathfrak h(A,B)$.  Its coefficient at the primary source $O$ is $+1$;
every other evaluation is at a singleton $u\in\{d_1,\ldots,d_\ell,v\}$.
Since $A\cap B=\varnothing$ on $D$,
\begin{equation}\label{eq:pair-test-shape}
 \mathfrak h(A,B)=b_O(A,B)+
   \sum_{u\in\{d_1,\ldots,d_\ell,v\}}\lambda_u\one\{u\in A\}
\end{equation}
for constants $\lambda_u$.  In particular, for fixed $A$ the map
$B\mapsto\mathfrak h(A,B)$ is antitone.  This observation is why the
argument works for one set-valued primary source but keeps all decoys and
the second observer singleton-valued.

Here is the exact reduction lemma.  Under the conditional law on $D$, let
$P_A$ and $P_B$ denote conditional expectation given $A$ and $B$,
respectively, and put $T=P_A P_B$ on functions of $A$.  For every bounded
$g=g(A)$, with $u=P_B g$, the elementary conditional-covariance identity is
\begin{multline}\label{eq:two-half-cov}
 \operatorname{Cov}_D(g,\mathfrak h)
 -\operatorname{Cov}_D(Tg,\mathfrak h)\\
 =\EE_D[\operatorname{Cov}(g,\mathfrak h\mid B)]
  +\EE_D[\operatorname{Cov}(u,\mathfrak h\mid A)].
\end{multline}
The conditional cluster decomposition and \eqref{eq:guard-W-def} give
the exact normalization
\[
 \EE_D[\operatorname{Cov}(g,\mathfrak h\mid B)]
 =\delta^{-1}\mathcal W_{x,Y,{\bf d}}^{O,v}(g).
\]
Here and below $\mathcal W(g)$ denotes \eqref{eq:guard-W-def} with the
cluster functional $\psi$ replaced by $g$.
Thus the first term is nonnegative.  The factor $\delta^{-1}$ is
essential: $\mathcal W$ is the unnormalised integral over $D$, whereas
\eqref{eq:two-half-cov} is written under the conditional law.  If $g$ is
increasing, then
$u(B)=\EE[g(A)\mid B]$ is antitone: conditionally on $B$, the cluster
$A$ is the $x$-cluster in the graph with every pair meeting $B$ deleted,
and the common-edge coupling for $B\subseteq B'$ makes the latter cluster
smaller.  This is the existing lemma \lean{BHK2006.condS\_antitone}.
By \eqref{eq:pair-test-shape}, $B\mapsto\mathfrak h(A,B)$ is also
antitone for fixed $A$.  Conditional on $A$, write
$B=\lean{BHK2006.halfT}(A,\eta)$, where $\eta$ is a fresh product
configuration in the graph with the pairs meeting $A$ deleted.  The map
$\eta\mapsto B$ is increasing.  Thus both
$u(B)$ and $\mathfrak h(A,B)$ are decreasing functions of $\eta$, and
\lean{BHK2006.harris\_anti\_anti} makes the second term in
\eqref{eq:two-half-cov} nonnegative.

For clarity, the other property needed in the iteration follows from the
same coupling.  If $A\subseteq A'$, then the conditional $Y$-cluster
sampled after deleting $A'$ is contained in the one sampled after
deleting $A$.  Since $u$ is antitone,
\[
 (Tg)(A)=\EE[u(B)\mid A]
 \leq \EE[u(B)\mid A']=(Tg)(A').
\]
Hence the BHK Gibbs operator $T$ preserves monotonicity; it plainly
preserves constants and bounds.  Under the present nondegeneracy
assumption it also has a uniform regeneration.  If $Y\ne\varnothing$,
close every pair meeting $Y$ in the $B$-resampling half-step.  This has
conditional probability at least
\[
 \epsilon:=\prod_{e:e\cap Y\ne\varnothing}(1-p_e)>0,
\]
independently of the old $A$, and forces $B=Y$.  The following
$A$-resampling then has a fixed law, independent of the old state.  Thus
the transition kernel is of the form
$\epsilon\nu+(1-\epsilon)Q_A$, and
\[
 \operatorname{osc}(T^n g)
 \leq(1-\epsilon)^n\operatorname{osc}(g).
\]
If $Y=\varnothing$, the first application of $T$ is already constant.
Since minima of $T^n g$ increase and maxima decrease, $T^n g$ converges
uniformly to a constant.  The same statements apply with $g$ replaced by
$T^n g$: the operator $T$ preserves nonnegativity and monotonicity, so the
first term in \eqref{eq:two-half-cov} is
$\delta^{-1}\mathcal W(T^n g)\geq0$, and the second is nonnegative by the
same antitone--antitone Harris argument.  Iterating
\eqref{eq:two-half-cov} consequently
gives
\[
 \operatorname{Cov}_D(g,\mathfrak h)
 \geq \operatorname{Cov}_D(T^n g,\mathfrak h)\longrightarrow0.
\]
Finally, direct expansion of \eqref{eq:guardGamma} and linearity of the
level forms give
\[
 L_\ell\Gamma_{x,Y}^{\psi}(O)
 -pL_\ell\Gamma_{x,Y}^{\psi}(\{v\})
 =\delta^2\operatorname{Cov}_D(g,\mathfrak h).
\]
Thus the last nonnegative covariance is exactly the guarded CSH margin
up to the positive factor $\delta^2$.  This proves the theorem.  Notice
that this
two-half-step calculation is essential: a fixed global kernel
$\alpha_{R,Y}^{V}$ does \emph{not} become $j_{\{x\}\mid Z}(R)$ merely by
conditioning on $C_Y=Z$.
\end{proof}

\begin{remark}[what is new, and what is reused]
The genuinely new declarations are the guarded event
\eqref{eq:guard-event}, the kernel \eqref{eq:alpha-kernel}, the exchange
lemma \eqref{eq:guard-exchange}, the additional $R\ncon N$ guard in
the generic two-source induction, the set-source Markov identity
\eqref{eq:guard-markov-a}, its nonnegative correction
\eqref{eq:guard-error}, and the two-half-step reduction
\eqref{eq:two-half-cov}.  The BHK four-product theorem,
set-cluster exploration, stopped Harris, Ahlswede--Daykin, the general
Gibbs reduction, and the residual functional are already proved in the
repository.  In particular, Theorem~\ref{thm:guardCSH} does not assume a
new correlation conjecture.  Omitting the guard would make the later
peeling step false; factoring the unevaluated source components would
incorrectly delete the positive error in \eqref{eq:guard-unfold}.
\end{remark}

\subsection{Set-source surplus transfer}

For $B\subseteq V$ write $D_B(k)=\{k\ncon B\}$.  If $k\notin Y\cup T$,
put
\begin{equation}\label{eq:EkJk}
 E_k:=D_{Y\cup T}(k),\qquad
 Q_{S,k}:=D_{Y\cup T}(S)\cap\{S\con k\}.
\end{equation}
For later use abbreviate the guarded one-decoy profile by
\begin{equation}\label{eq:c-general}
 c_{z,B}(R):=
 \frac{\PP(\mathsf H_{\{z\}\mid B}(R))}{\PP(z\ncon B)}.
\end{equation}
For a relay set $T$ and an additional decoy list
${\bf d}=(d_1,\ldots,d_\ell)$, set
\[
 B_j=Y\cup T\cup\{d_i:i<j\},\qquad
 B_*=Y\cup T\cup\{d_1,\ldots,d_\ell\}.
\]
Use the profiles $c_{d_j,B_j}$ in the recursion \eqref{eq:guard-slack}
and the price
\[
 p_{O,v}^{Y,T,{\bf d}}
 =\frac{\PP(\mathsf H_{\{v\}\mid B_*}(O))}
        {\PP(v\ncon B_*)}.
\]
We write
\begin{equation}\label{eq:surplus-general-margin}
 \operatorname{Marg}_{Y,T,{\bf d}}^{O,v}[h]
 :=L_\ell h(O)-p_{O,v}^{Y,T,{\bf d}}L_\ell h(\{v\}).
\end{equation}
This definition removes a possible ambiguity below: the base avoided set
for a surplus margin is $Y\cup T$, not merely $Y$.

\begin{theorem}[set-source avoided surplus transfer]\label{thm:setS5}
Let $\Phi$ be increasing and let $r$ be compatible with the conditional
means \eqref{eq:condmean}.  Suppose that $S,Y,T,\{k\}$ are pairwise
disjoint.  Then, for arbitrary edge weights in $[0,1]$,
\begin{equation}\label{eq:setS5}
 \PP(E_k)\,\Sur_{S,Y}(T;r,\Phi)
 \ \geq\
 \PP(Q_{S,k})\,\Sur_{k,Y}(T;r,\Phi).
\end{equation}
\end{theorem}

\begin{proof}
First take nondegenerate weights and decoys disjoint from all displayed
sets.  We prove, simultaneously for every such decoy list ${\bf d}$ and
singleton $v$, the stronger statement
\begin{equation}\label{eq:setS5D}
 \operatorname{Marg}_{Y,T,{\bf d}}^{S,v}
       [R\mapsto\Sur_{R,Y}(T;r,\Phi)]\geq0.
\end{equation}
Induct on $|T|$.

If $T=\varnothing$, then $R\con T$ is impossible and
$\Sur_{R,Y}(T)=0$ for every source $R$.  Thus
\eqref{eq:setS5D}, and hence \eqref{eq:setS5}, holds with equality.
Assume from now on that $T\ne\varnothing$.

Let $z$ have maximal rank and write $T'=T\setminus\{z\}$.  Set
\[
 E_z=D_{Y\cup T'}(z),\quad d_z=\PP(E_z),\quad
 \kappa_z=m_z^Y d_z-\EE[\Phi(C_z);E_z]
\]
and
\begin{equation}\label{eq:theta-z}
 \Theta_z(R)=\EE[(\Phi(C_z)-m_z^Y);
                    \mathsf H_{\{z\}\mid Y\cup T'}(R)].
\end{equation}
For a singleton $R=\{u\}$, the usual first-pattern partition gives the
exact peeling identity
\begin{equation}\label{eq:singleton-peel-Y}
 \Sur_{u,Y}(T)=\Sur_{u,Y}(T')+\Theta_z(\{u\}).
\end{equation}
For $R=S$, the same partition leaves
$\Phi(C_S)$ in the new term, on the \emph{same guarded event} appearing
in \eqref{eq:theta-z}.  On that event $S\con z$, hence
$C_z\subseteq C_S$; monotonicity gives the one-sided inequality
\begin{equation}\label{eq:set-peel-Y}
 \Sur_{S,Y}(T)\geq\Sur_{S,Y}(T')+\Theta_z(S).
\end{equation}
This is the only point at which a source union differs from a singleton.
In the level margin the value at $S$ occurs with coefficient $+1$ and
all other evaluations are the singleton decoys and $v$.  Therefore
\eqref{eq:set-peel-Y} and \eqref{eq:singleton-peel-Y} preserve the
inequality after applying the level form.

For every finite source $R$, direct expansion of
\eqref{eq:guardGamma}
gives
\begin{equation}\label{eq:theta-cov}
 d_z\Theta_z(R)=
 \Gamma_{z,Y\cup T'}^{\Phi}(R)
 -\kappa_z d_z c_{z,Y\cup T'}(R).
\end{equation}
The guarded-source CSH theorem makes the margin of the first term
nonnegative.  Apply it once more to the increasing vertex-cluster
functional
\[
 \psi_{{\rm iso},z}(K)=\one\{K\setminus\{z\}\ne\varnothing\}.
\]
Every evaluated source is distinct from $z$.  If
$D=\{z\ncon Y\cup T'\}$ and
$\rho_z=\PP(D,C_z=\{z\})$, direct expansion gives
\begin{equation}\label{eq:iso-guard-identity}
 \Gamma_{z,Y\cup T'}^{\psi_{{\rm iso},z}}(R)
 =\rho_z\,\PP(\mathsf H_{\{z\}\mid Y\cup T'}(R))
 =\rho_z d_z c_{z,Y\cup T'}(R).
\end{equation}
Under nondegenerate weights $\rho_z d_z>0$.  Guarded CSH for
$\psi_{{\rm iso},z}$ therefore says precisely that the required margin of
$c_{z,Y\cup T'}$ is nonnegative.  Using a vertex-cluster indicator here
of the event $K\ne\varnothing$ would be wrong, because a vertex cluster
is never empty; the isolated state is $C_z=\{z\}$.

Finally, centering over $D_Y(z)$ and partitioning according as $z$ hits
$T'$ or avoids it gives
\begin{equation}\label{eq:kappasur-transfer}
 \Sur_{z,Y}(T')-\kappa_z
 =\sum_{x\in T'}\PP(J_x^{\{z\},Y})(m_z^Y-m_x^Y)\geq0.
\end{equation}
The last inequality is exactly maximality of the rank of $z$.  Prepending
$z$ to the decoy list gives the algebraic identity
\[
 \operatorname{Marg}_{z::{\bf d}}[\Sur(T')]
 =\operatorname{Marg}_{\bf d}[\Sur(T')]
  -\Sur_{z,Y}(T')\operatorname{Marg}_{\bf d}[c_z].
\]
Equations \eqref{eq:set-peel-Y}--\eqref{eq:kappasur-transfer}, the two
nonnegative CSH margins, and the induction hypothesis therefore give
\[
 d_z\operatorname{Marg}_{\bf d}[\Sur(T)]
 \geq d_z\operatorname{Marg}_{z::{\bf d}}[\Sur(T')]
 \geq0.
\]
At decoy level zero, multiply out the observer price.  With $v=k$ this
is precisely \eqref{eq:setS5}.

It remains to remove the nondegeneracy restriction.  This closure is
proved explicitly just below; it applies to the cross-multiplied
inequality \eqref{eq:setS5} and completes the proof.
\end{proof}

\subsection{Ranks, zero denominators, and boundary weights}

The closure used above is short but essential.  Call $x\in T$ active if
$d_Y(x)>0$.  If $x$ is inactive, every event
$D_Y(S)\cap\{S\con x\}$ is null, because it is contained in
$D_Y(x)$.  Inactive relays can therefore be deleted without changing any
integral or first-pattern term.  For the active relays, every compatible
rank gives the rank-free identity
\begin{equation}\label{eq:rank-free-surplus}
 \Sur_{S,Y}(T;r,\Phi)=
 \EE\!\left[
  \one\{D_Y(S),S\con T\}
  \left(\Phi(C_S)-
   \min_{x\in T\cap C_S}m_x^Y\right)\right].
\end{equation}
Indeed, the first relay in $C_S$ has the least conditional mean among
$T\cap C_S$; ties have the same value.  Formula
\eqref{eq:rank-free-surplus} also proves independence from tie-breaking.

Approximate an arbitrary weight vector $w\in[0,1]^E$ by nondegenerate
vectors $w_n\in(0,1)^E$.  There are only finitely many relay orders, so
pass to a subsequence on which the compatible order is fixed.  For every
relay active at $w$, its avoidance probability stays positive and its
conditional mean converges to the mean at $w$.  Terms involving a relay
inactive at $w$ are supported on $D_Y(x)$ and hence tend to zero.  Thus
both surpluses in \eqref{eq:setS5} converge to the active-relay
expressions at $w$.  The two probability prefactors are polynomials in
the edge weights.  Passing to the limit proves \eqref{eq:setS5} at $w$.
This argument neither chooses a limiting minimizer nor divides by a
quantity that may vanish.

\subsection{The conditional top-relay estimate}

The following elementary lemma is useful both for checking the signs and
for formalization.

\begin{lemma}[set-source top-relay estimate]\label{lem:toprelay}
Let $T_0\subseteq V\setminus Y$, let $k\notin Y\cup T_0$, and set
\[
 E=D_{Y\cup T_0}(k),\qquad
 J=D_{Y\cup T_0}(S)\cap\{S\con k\},\qquad
 \kappa=m_k^Y\PP(E)-\EE[\Phi(C_k);E].
\]
Then
\begin{equation}\label{eq:toprelay}
 \PP(E)\,
 \EE[(\Phi(C_S)-m_k^Y);J]
 \geq-\PP(J)\kappa.
\end{equation}
\end{lemma}

\begin{proof}
On $J$ one has $C_k\subseteq C_S$, hence it is enough to replace
$\Phi(C_S)$ by $\Phi(C_k)$ on the left.

We claim
\begin{equation}\label{eq:conditional-assoc-source}
 \PP(E)\EE[\Phi(C_k);J]
 \geq \PP(J)\EE[\Phi(C_k);E].
\end{equation}
To see this directly, condition on $E$ and then on $C_k=K$.  Here
$K\cap(Y\cup T_0)=\varnothing$.  Given $K$, the unexplored edges off $K$
are fresh, and the conditional probability of $J$ is
\[
 q(K)=\one\{K\cap S\ne\varnothing\}
 \PP_{G\setminus K}\bigl(C_{S\setminus K}\cap(Y\cup T_0)=\varnothing\bigr),
\]
where $G\setminus K$ means that all pairs meeting $K$ are deleted.  If
$K\subseteq K'$ and both avoid $Y\cup T_0$, then $q(K)\leq q(K')$:
the first indicator can only increase, the residual source shrinks, and
the residual graph loses edges.  Thus both $K\mapsto\Phi(K)$ and $q(K)$
are increasing.  Conditional positive association for $C_k$ given
$k\ncon Y\cup T_0$ gives \eqref{eq:conditional-assoc-source}.

Subtracting $m_k^Y\PP(E)\PP(J)$ from the two sides of
\eqref{eq:conditional-assoc-source} gives \eqref{eq:toprelay}.
\end{proof}

\subsection{Source-set avoided GEN}

The transfer theorem was stated for disjoint sources and relays.  The
following one-cluster estimate handles all overlaps and is also useful as
a separate formalization target.

\begin{lemma}[overlapping source relay]\label{lem:source-overlap}
Let $t\in S\setminus Y$ and let $\Phi$ be increasing.  Then
\begin{equation}\label{eq:source-overlap}
 d_Y(t)\,\EE[\Phi(C_S);D_Y(S)]
 \geq
 \PP(D_Y(S))\,\EE[\Phi(C_t);D_Y(t)].
\end{equation}
\end{lemma}

\begin{proof}
If $d_Y(t)=0$, then $D_Y(S)\subseteq D_Y(t)$ is null and both sides
vanish.  Suppose it is positive.  For a possible cluster $K=C_t$ with
$K\cap Y=\varnothing$, put
\[
 \alpha(K)=\PP_{G\setminus K}
  \bigl(C_{S\setminus K}\cap Y=\varnothing\bigr).
\]
If $K\subseteq K'$, the residual source shrinks and the residual graph
loses pairs, so $\alpha(K)\leq\alpha(K')$.  Exploration of $C_t$ gives
the two exact tower identities
\begin{align*}
 \EE[\alpha(C_t);D_Y(t)]&=\PP(D_Y(S)),\\
 \EE[\Phi(C_t)\alpha(C_t);D_Y(t)]
   &=\EE[\Phi(C_t);D_Y(S)].
\end{align*}
Conditional positive association of $C_t$ given $D_Y(t)$, applied to
the increasing functions $K\mapsto\Phi(K)$ and $K\mapsto\alpha(K)$,
therefore yields
\[
 d_Y(t)\EE[\Phi(C_t);D_Y(S)]
 \geq
 \EE[\Phi(C_t);D_Y(t)]\PP(D_Y(S)).
\]
The association statement is unchanged after adding a constant to
$\Phi$, so no sign assumption is needed.  Finally
$C_t\subseteq C_S$ pointwise, and monotonicity of $\Phi$ proves
\eqref{eq:source-overlap}.
\end{proof}

\begin{theorem}[source-set avoided first-relay inequality]\label{thm:SAG}
Let $S,Y,T\subseteq V$, with $T\cap Y=\varnothing$.  Let $\Phi:2^V\to
\mathbb R$ be increasing, and rank $T$ compatibly with the conditional
means $m_x^Y$.  Then
\begin{equation}\label{eq:SAG}
 \boxed{\quad
 \sum_{x\in T}\PP(J_x^{S,Y})m_x^Y
 \leq
 \EE[\Phi(C_S);D_Y(S),S\con T].
 \quad}
\end{equation}
Equivalently, $\Sur_{S,Y}(T;r,\Phi)\geq0$.
\end{theorem}

\begin{proof}
If $S\cap Y\ne\varnothing$, then $D_Y(S)$ is empty and both sides of
\eqref{eq:SAG} vanish.  Suppose next that $t\in S\cap T$.  If $t$ is
inactive, then $D_Y(S)\subseteq D_Y(t)$ is null.  Otherwise $S\con T$
is automatic, and on every first-relay pattern the relay met first has
conditional mean at most $m_t^Y$, since $t\in C_S$ is itself available.
Hence Lemma~\ref{lem:source-overlap} gives
\[
 \sum_{x\in T}\PP(J_x^{S,Y})m_x^Y
 \leq \PP(D_Y(S))m_t^Y
 \leq \EE[\Phi(C_S);D_Y(S)],
\]
which is \eqref{eq:SAG}.  We may therefore assume for the rest of the
proof that $S,Y,T$ are pairwise disjoint; this is exactly the domain in
which Theorem~\ref{thm:setS5} is used.

Adding a constant to $\Phi$ adds that constant times the two equal
quantities in \eqref{eq:partition}; hence we may assume $\Phi\geq0$.
The empty relay case is immediate.  Suppose $T\ne\varnothing$.

Let $k$ be the rank-maximal relay and put $T_0=T\setminus\{k\}$.  Set
\[
 E=D_{Y\cup T_0}(k),\qquad
 J=D_{Y\cup T_0}(S)\cap\{S\con k\},
\]
and define $\kappa$ as in Lemma~\ref{lem:toprelay}.  The pattern $J$ is
exactly the new first-relay pattern for $k$, and therefore
\begin{equation}\label{eq:surplus-peel}
 \Sur_{S,Y}(T;r,\Phi)
 =\Sur_{S,Y}(T_0;r,\Phi)
  +\EE[(\Phi(C_S)-m_k^Y);J].
\end{equation}

The centered integral of $\Phi(C_k)$ over $D_Y(k)$ is zero.  Splitting
$D_Y(k)$ according as $k$ hits $T_0$ or avoids it gives
\begin{align}
 \kappa
 &=\EE[(\Phi(C_k)-m_k^Y);D_Y(k),k\con T_0],\label{eq:kappa1}\\
 \Sur_{k,Y}(T_0;r,\Phi)-\kappa
 &=\sum_{x\in T_0}\PP(J_x^{\{k\},Y})
       (m_k^Y-m_x^Y)\geq0.\label{eq:kappa2}
\end{align}
The last inequality uses maximality of the rank of $k$.

Multiply \eqref{eq:surplus-peel} by $\PP(E)$.  Apply
Theorem~\ref{thm:setS5} to the first term and
Lemma~\ref{lem:toprelay} to the second.  We obtain
\[
 \PP(E)\Sur_{S,Y}(T;r,\Phi)
 \geq \PP(J)\bigl(\Sur_{k,Y}(T_0;r,\Phi)-\kappa\bigr)\geq0
\]
by \eqref{eq:kappa2}.  In the nondegenerate case $\PP(E)>0$, so the
desired surplus is nonnegative.  The denominator-free boundary argument
and rank-free continuity argument preceding this subsection pass the
inequality to arbitrary weights.
\end{proof}

\section{The monotone projection}\label{sec:projection}

We now turn Theorem~\ref{thm:SAG} into a comparison with a fixed relay.
Fix $a\in V$.  For $K\subseteq V$ with $a\notin K$, let $G\setminus K$
denote the random graph obtained by deleting all pairs with an endpoint in
$K$, and define
\begin{equation}\label{eq:projection}
 \phi_a(K):=F(K)-\EE F(C_a^{G\setminus K}).
\end{equation}

\begin{lemma}[projection facts]\label{lem:projection}
Suppose $F$ is increasing.
\begin{enumerate}[label=\textup{(P\arabic*)},leftmargin=3em]
\item The map $K\mapsto\phi_a(K)$ is increasing on
  $\{K:a\notin K\}$ and has an increasing extension to $2^V$.
\item For every source set $S$ and every event $H$ determined by $C_S$,
  \begin{equation}\label{eq:projection-source}
   \EE[(F(C_S)-F(C_a));S\ncon a,H]
   =\EE[\phi_a(C_S);S\ncon a,H].
  \end{equation}
\item For every vertex $x$,
  \begin{equation}\label{eq:projection-mean}
   \EE[\phi_a(C_x);x\ncon a]
   =\EE F(C_x)-\EE F(C_a).
  \end{equation}
\end{enumerate}
\end{lemma}

\begin{proof}
If $K\subseteq L$ and neither contains $a$, then $F(K)\leq F(L)$, while
deleting all pairs meeting $L$ leaves no more open paths from $a$ than
deleting all pairs meeting $K$.  Thus the second term in
\eqref{eq:projection} decreases, proving (P1) on the indicated subposet.
Give every set containing $a$ the value
$\max_{a\notin K}\phi_a(K)$; this is an increasing extension.  A final
constant shift makes it nonnegative if desired.

On $S\ncon a$, condition on $C_S=K$.  Every pair from $K$ to its
complement is closed, and the edges wholly outside $K$ retain their
independent laws.  Consequently the conditional law of $C_a$ is exactly
the law of $C_a^{G\setminus K}$.  This proves (P2).  Apply (P2) with
$S=\{x\}$ and $H$ the sure event.  On $x\con a$ the two clusters, and
hence their $F$-values, agree.  Adding this zero contribution proves
(P3).
\end{proof}

The repository already contains the singleton versions:
\begin{itemize}[nosep]
\item \lean{CovTau.monotone\_projFun};
\item \lean{CovTau.tower\_clusterFun};
\item \lean{CovTau.setIntegral\_sub\_eq\_projFun}.
\end{itemize}
The proof above shows that the same interface works for the explored
cluster of a source set.

\section{A fixed minimizer works: Conjecture 4}\label{sec:C4}

\begin{theorem}[fixed-minimizer source-set comparison]\label{thm:fixedmin}
Let $F:2^V\to\mathbb R$ be increasing, let $A\ne\varnothing$, and choose
$a\in A$ such that
\begin{equation}\label{eq:amin}
 \EE F(C_a)=\min_{x\in A}\EE F(C_x).
\end{equation}
Then, for every finite source set $S$,
\begin{equation}\label{eq:fixedavoid}
 \EE[(F(C_S)-F(C_a));S\ncon a,\ S\con A\setminus\{a\}]
 \geq0.
\end{equation}
Consequently
\begin{equation}\label{eq:fixedfull}
 \EE[F(C_S);S\con A]\geq\EE[F(C_a);S\con A].
\end{equation}
For $S=\{o\}$, the contribution on $o\con a$ vanishes identically, so
\eqref{eq:fixedfull} holds with equality removed only on the avoided part.
\end{theorem}

\begin{proof}
Put $Y=\{a\}$ and $T=A\setminus\{a\}$.  Use the increasing projection
$\phi_a$ from Lemma~\ref{lem:projection}.  By
\eqref{eq:projection-mean} and \eqref{eq:amin}, every conditional relay
mean is nonnegative:
\begin{equation}\label{eq:positive-projmean}
 \EE[\phi_a(C_x)\mid x\ncon a]
 =\frac{\EE F(C_x)-\EE F(C_a)}{\PP(x\ncon a)}\geq0.
\end{equation}
Relays for which the denominator is zero never occur in an $a$-avoiding
first pattern and may be removed.  Rank the others by the quantities in
\eqref{eq:positive-projmean}.  Theorem~\ref{thm:SAG}, followed by
Lemma~\ref{lem:projection}\textup{(P2)}, gives
\begin{align*}
 &\EE[(F(C_S)-F(C_a));S\ncon a,S\con T]\\
 &\quad=\EE[\phi_a(C_S);S\ncon a,S\con T]\\
 &\quad\geq
 \sum_{x\in T}\PP(J_x^{S,\{a\}})
       \EE[\phi_a(C_x)\mid x\ncon a]\geq0.
\end{align*}
This is \eqref{eq:fixedavoid}.

On the remaining part of $\{S\con A\}$, the source cluster union meets
$a$.  Hence $C_a\subseteq C_S$ and $F(C_a)\leq F(C_S)$.  Adding that
nonnegative contribution proves \eqref{eq:fixedfull}.
\end{proof}

\begin{corollary}[Conjecture 4, in a stronger form]\label{cor:C4}
For a singleton source $S=\{o\}$, Theorem~\ref{thm:fixedmin} says that the
particular unconditional minimizer $a$ in \eqref{eq:amin} satisfies
\begin{equation}\label{eq:C4fixed}
 \EE[F(C_o);o\con A]\geq\EE[F(C_a);o\con A].
\end{equation}
Thus Conjecture~4 holds.  This also answers the first \KN\ Question~7:
the minimizer of the post-FKG quantity is a valid relay in the pre-FKG
comparison.
\end{corollary}

\section{Conjectures 2, 1, and 3}\label{sec:C123}

\begin{corollary}[strong Conjecture 2]\label{cor:C2}
For all $o,b,A$,
\begin{equation}\label{eq:C2strong-proof}
 \PP(o\con b,o\con A)
 \geq \min_{x\in A}\PP(x\con b,o\con A).
\end{equation}
In particular, Conjecture~2 holds.
\end{corollary}

\begin{proof}
Apply Corollary~\ref{cor:C4} to
$F(K)=\one\{b\in K\}$.  Then $F(C_x)=\one\{x\con b\}$, and
\eqref{eq:C4fixed} is at least as strong as
\eqref{eq:C2strong-proof}.  Finally
$\PP(o\con b)\geq\PP(o\con b,o\con A)$.
\end{proof}

\begin{corollary}[Conjecture 1]\label{cor:C1}
The multiplicative inequality \eqref{eq:C1} holds for every finite
weighted graph.
\end{corollary}

\begin{proof}
Both $\{o\con A\}$ and $\{x\con b\}$ are increasing events.  Harris--FKG
and Corollary~\ref{cor:C2} give
\begin{align*}
 \PP(o\con b)
 &\geq\min_{x\in A}\PP(o\con A,x\con b)\\
 &\geq\min_{x\in A}\PP(o\con A)\PP(x\con b)\\
 &=\PP(o\con A)\min_{x\in A}\PP(x\con b).
\end{align*}
\end{proof}

\begin{remark}
Conjecture~1 also follows directly from the repository's unavoided GEN
theorem by taking $F(K)=\one\{b\in K\}$ and ordering relays by
$\PP(x\con b)$.  The route above is included to display the implication
chain in the Kozma--Nitzan paper.
\end{remark}

\begin{corollary}[Conjecture 3]\label{cor:C3}
Conjecture~3 holds; for $0<\varepsilon\leq1$ one may take
$\delta=\varepsilon/2$.
\end{corollary}

\begin{proof}
Under its hypotheses, Corollary~\ref{cor:C1} gives
\[
 \PP(o\con b)>(1-\delta)^2
 =1-\varepsilon+\frac{\varepsilon^2}{4}>1-\varepsilon.
\]
For $\varepsilon>1$ the assertion is trivial, for example with
$\delta=1/2$.
\end{proof}

\begin{proposition}[countable-graph closure]\label{prop:countable}
On a countable graph with independent, possibly inhomogeneous bond
weights, Conjecture~1 holds with $\inf_{x\in A}$ in place of the minimum
when $A$ is infinite.  Consequently Conjecture~3 holds there with the
same graph-independent choice $\delta=\varepsilon/2$.
\end{proposition}

\begin{proof}
First let $A$ be finite.  Exhaust the countable edge set by finite edge
sets, retaining their original weights and declaring every other edge
closed.  Each connection event between specified vertices increases to
the corresponding event in the full graph, because every connecting path
is finite.  Apply Corollary~\ref{cor:C1} at every finite stage and use
continuity from below; the minimum of finitely many convergent sequences
converges to the minimum of their limits.

For countably infinite $A$, exhaust it by nonempty finite sets $A_m$.
Then $\{o\con A_m\}$ increases to $\{o\con A\}$ and
$\min_{x\in A_m}\PP(x\con b)$ decreases to
$\inf_{x\in A}\PP(x\con b)$.  Passing to the limit proves the first
claim.  Under the hypotheses of Conjecture~3, continuity from below also
provides a finite $A_0\subseteq A$ with
$\PP(o\con A_0)>1-\delta$.  The already proved finite-relay inequality
then gives the same estimate as in Corollary~\ref{cor:C3}.
\end{proof}

\begin{corollary}[the main critical-percolation consequence]
For Bernoulli bond percolation on $\mathbb Z^d$, $d\geq2$,
\[
 \PP_{p_c}(|C(0)|=\infty)=0.
\]
\end{corollary}

\begin{proof}
Proposition~\ref{prop:countable} supplies Conjecture~3 in the uniform form
assumed by Theorem~6 of Kozma--Nitzan.  That theorem gives the displayed
conclusion.
\end{proof}

\section{The new edge argument: Conjecture 6}\label{sec:C6}

The source-set extension makes the edge-contraction conjecture a short
corollary.  We isolate the deterministic observation on which it rests.

\begin{lemma}[an internal edge does not change a source union]\label{lem:internal}
Let $e=\{v,w\}$, let $S=\{v,w\}$, and couple all weightings of $e$ by
using the same states for every other pair.  Then the following objects do
not depend on the state of $e$:
\begin{enumerate}[label=\textup{(E\arabic*)},leftmargin=3em]
\item the vertex set $C_S=C_v\cup C_w$;
\item every event $\{S\con B\}$ and $\{S\ncon B\}$;
\item on $S\ncon a$, the cluster $C_a$.
\end{enumerate}
When $e$ is forced open, $C_v=C_w=C_S$.
\end{lemma}

\begin{proof}
Opening an edge with both endpoints in the source set merely joins two
components already included in their union; it adds no vertex to that
union.  This proves (E1), and (E2) follows.  On $S\ncon a$, the cluster of
$a$ contains neither endpoint of $e$, so changing $e$ cannot change that
cluster.  The last assertion is immediate when $e$ is open.
\end{proof}

\begin{theorem}[fixed-minimizer contraction; stronger than Conjecture 6]
\label{thm:C6strong}
Let $e=\{v,w\}$ and let $a\in A$ minimize $\PP_G(x\con b)$ over $A$.
Then, without any endpoint hypotheses,
\begin{equation}\label{eq:pinned-preFKG}
 \PP_{G/e}(v\con b,v\con A)
 \geq \PP_{G/e}(a\con b,v\con A),
\end{equation}
and consequently
\begin{equation}\label{eq:C6proved}
 \PP_{G/e}(v\con b)
 \geq\PP_{G/e}(v\con A)\PP_{G/e}(a\con b).
\end{equation}
\end{theorem}

\begin{proof}
Use $F(K)=\one\{b\in K\}$ in Theorem~\ref{thm:fixedmin}, in the original
weighted graph $G$, with the source set $S=\{v,w\}$.  It gives
\begin{equation}\label{eq:source-edge-positive}
 \EE_G[(F(C_S)-F(C_a));S\ncon a,S\con A\setminus\{a\}]
 \geq0.
\end{equation}
By Lemma~\ref{lem:internal}, the random variable and event in
\eqref{eq:source-edge-positive} are independent of the state of $e$.
We may therefore evaluate them with $e$ forced open.  In $G/e$, the
source union $C_S$ is $C_v$.  Moreover, on $v\con a$ the two indicators
$F(C_v)$ and $F(C_a)$ agree, while on $v\ncon a$ the event $v\con A$ is
the event $v\con A\setminus\{a\}$.  It follows that the left side of
\eqref{eq:source-edge-positive} is exactly
\[
 \PP_{G/e}(v\con b,v\con A)
 -\PP_{G/e}(a\con b,v\con A).
\]
This proves \eqref{eq:pinned-preFKG}.

Finally, Harris--FKG in $G/e$ gives
\[
 \PP_{G/e}(a\con b,v\con A)
 \geq\PP_{G/e}(a\con b)\PP_{G/e}(v\con A).
\]
Together with
$\PP_{G/e}(v\con b)\geq\PP_{G/e}(v\con b,v\con A)$ and
\eqref{eq:pinned-preFKG}, this proves \eqref{eq:C6proved}.
\end{proof}

\begin{corollary}[Conjecture 6]\label{cor:C6}
Conjecture~6 holds.  Its two assumptions concerning $v$ and $w$ are
superfluous.
\end{corollary}

\section{The numbered questions: common coefficients and candidate relays}
\label{sec:questions}

We now dispose of Questions~5 and 7--9.  The coefficient conclusion below
is stronger than the one printed in the paper.

\begin{theorem}[common coefficients; Question 5]\label{thm:Q5}
For fixed $G,o,A$ there are coefficients $c_x\geq0$, $x\in A$, with
$\sum_{x\in A}c_x=1$, such that simultaneously for every $b\in V$,
\begin{equation}\label{eq:Q5strong}
 \sum_{x\in A}c_x\PP(o\con A,x\con b)
 \leq \PP(o\con b,o\con A)
 \leq \PP(o\con b).
\end{equation}
In particular, Question~5 has an affirmative answer.
\end{theorem}

\begin{proof}
Let $E=\{o\con A\}$ and regard the following as vectors indexed by
$b\in V$:
\[
 r_b=\PP(E,o\con b),\qquad q_b^x=\PP(E,x\con b)\quad(x\in A).
\]
For every vector $\lambda=(\lambda_b)_{b\in V}\geq0$, the functional
\[
 F_\lambda(K)=\sum_{b\in K}\lambda_b
\]
is nonnegative and increasing.  Conjecture~4, already proved in
Corollary~\ref{cor:C4}, gives
\begin{equation}\label{eq:Q5-separation-test}
 \lambda\mathbin{\cdot}r
 =\EE[F_\lambda(C_o);E]
 \geq\min_{x\in A}\EE[F_\lambda(C_x);E]
 =\min_{x\in A}\lambda\mathbin{\cdot}q^x.
\end{equation}

Put $K=\operatorname{conv}\{q^x:x\in A\}\subseteq\mathbb R^V$ and
$D=r-\mathbb R_{\geq0}^{V}$.  Suppose $K\cap D=\varnothing$.  Strict
separation of the compact convex set $K$ from the closed convex set $D$
then gives a nonzero vector $\lambda$ such that
\[
 \inf_{z\in K}\lambda\mathbin{\cdot}z
 >\sup_{y\in D}\lambda\mathbin{\cdot}y.
\]
Since $D$ is unbounded in every negative coordinate direction, finiteness
of the supremum forces $\lambda\geq0$, and that supremum is
$\lambda\mathbin{\cdot}r$.  The infimum over $K$ is
$\min_{x\in A}\lambda\mathbin{\cdot}q^x$, contradicting
\eqref{eq:Q5-separation-test}.  Hence $K\cap D\ne\varnothing$.
Write a point in the intersection as
$\sum_{x\in A}c_x q^x=r-u$, where $c_x\geq0$, $\sum_x c_x=1$, and
$u\geq0$.  Coordinatewise this is the first inequality in
\eqref{eq:Q5strong}; the second is event inclusion.  The single vector
$(c_x)$ was chosen using all coordinates at once and is therefore
independent of $b$.
\end{proof}

\begin{corollary}[Questions 7 and 8]\label{cor:Q78}
Questions~7 and 8 both have affirmative answers, including every choice
among tied minimizers.
\end{corollary}

\begin{proof}
For Question~7, take $F(K)=\one\{b\in K\}$ in
Corollary~\ref{cor:C4}: the specified minimizer of
$\PP(x\con b)$ satisfies \eqref{eq:Qcandidate}.  For Question~8, the
specified $a$ realizes
$\min_{x\in A}\PP(o\con A,x\con b)$, so the strong form
\eqref{eq:C2strong-proof} of Conjecture~2 is exactly
\eqref{eq:Qcandidate}.
\end{proof}

\begin{theorem}[deleted-star minimizer; Question 9]\label{thm:Q9}
Let $H$ be obtained from $G$ by deleting every pair incident to $o$, and
let $a\in A$ minimize $\PP_H(x\con b)$ over $x\in A$.  Then
\eqref{eq:Qcandidate} holds in $G$.  Thus Question~9 has an affirmative
answer, again for every tied minimizer.
\end{theorem}

\begin{proof}
Expose all pairs incident to $o$.  For a fixed star configuration $\eta$,
let
\[
 N_\eta=\{u\ne o:\{o,u\}\text{ is open in }\eta\},
 \qquad S_\eta=\{o\}\cup N_\eta.
\]
The unexposed pairs have exactly the percolation law of $H$, and,
pointwise in their configuration,
\begin{equation}\label{eq:star-source-identities}
 C_o^G=C_{S_\eta}^{H},\qquad
 \{o\con_G A\}=\{S_\eta\con_H A\}.
\end{equation}
The inclusion of the isolated vertex $o$ in $S_\eta$ makes these formulas
valid without exceptions when $o\in A$ or $b=o$.

Put $F_b(K)=\one\{b\in K\}$.  On
$S_\eta\con_H a$, the full $G$-clusters of $o$ and $a$ coincide, so the
two integrands in \eqref{eq:Qcandidate} agree.  On
$S_\eta\ncon_H a$, no open star pair can change $C_a$, whence
$C_a^G=C_a^H$; moreover, on this event a hit of $A$ is a hit of
$A\setminus\{a\}$.  Therefore the difference between the two sides of
\eqref{eq:Qcandidate} has the exact decomposition
\begin{multline}\label{eq:Q9-decomposition}
 \PP_G(o\con b,o\con A)-\PP_G(a\con b,o\con A)\\
 =\sum_\eta\PP(\eta)\,
 \EE_H\!\left[(F_b(C_{S_\eta})-F_b(C_a));
 S_\eta\ncon a,\ S_\eta\con A\setminus\{a\}\right].
\end{multline}
The minimizer assumption is
$\EE_H F_b(C_a)\leq\EE_H F_b(C_x)$ for every $x\in A$.  Apply
Theorem~\ref{thm:fixedmin} in the graph $H$ to each deterministic source
set $S_\eta$.  Every expectation on the right of
\eqref{eq:Q9-decomposition} is nonnegative.  Averaging proves the theorem.
This also handles source--relay overlaps and star configurations of zero
probability; no division is used.
\end{proof}

\section{Dependency graph and audit checklist}\label{sec:audit}

The logical chain is
\[
 \boxed{\mathrm{CSH}_{\mathrm{set\ source}}}
 \Longrightarrow
 \boxed{\mathrm{SAG}}
 \Longrightarrow
 \boxed{\text{fixed-minimizer C4}}
 \Longrightarrow
 \boxed{\mathrm{C4}\Rightarrow\mathrm{C2}\Rightarrow\mathrm{C1}
        \Rightarrow\mathrm{C3}},
\]
and independently
\[
 \boxed{\text{fixed-minimizer C4 for }S=\{v,w\}}
 \Longrightarrow
 \boxed{\text{fixed-minimizer pre-FKG in }G/e}
 \Longrightarrow
 \boxed{\mathrm{C6}}.
\]
The remaining numbered questions branch off as follows:
\[
 \boxed{\mathrm{C4}}\Longrightarrow\boxed{\mathrm{Q5}},\qquad
 \boxed{\text{fixed-minimizer C4}}\Longrightarrow\boxed{\mathrm{Q7}},
 \qquad
 \boxed{\text{strong C2}}\Longrightarrow\boxed{\mathrm{Q8}},
\]
and
\[
 \boxed{\text{fixed-minimizer C4 for arbitrary }S}
 \mathrel{+}\boxed{\text{deleted-star conditioning}}
 \Longrightarrow\boxed{\mathrm{Q9}}.
\]

For auditing, the sign-sensitive steps are only these:
\begin{enumerate}[leftmargin=2.2em]
\item The set-source Markov exploration gives
  $-\widehat\Phi(C_R)$, not $-\widehat\Phi(C_d)$.  Their difference is
  the nonnegative error \eqref{eq:guard-error}; because the level split
  subtracts the centred bracket, that error enters
  \eqref{eq:guard-unfold} with a plus sign.
\item Every evaluation made by a tail level form or the final observer is
  a singleton.  Hence \eqref{eq:guard-error-singleton}, rather than a
  positivity-preservation assertion for a signed level operator, is what
  preserves the preceding inequality.
\item In the nonempty-intersection proof of \eqref{eq:guard-A2}, the
  cross-deleted sets $P,P'$ are essential: they make the arguments exactly
  $s\cup t$ and $s\cap t$ before applying four-functions.
\item In \eqref{eq:kappa2}, the peeled relay is rank \emph{maximal}, so
  $m_k^Y-m_x^Y\geq0$.
\item In Lemma~\ref{lem:toprelay}, $C_k\subseteq C_S$ and $\Phi$ is
  increasing, so replacing $\Phi(C_S)$ by $\Phi(C_k)$ is a lower bound.
\item The residual probability $q(K)$ in that lemma is increasing: a
  larger explored cluster deletes more possible paths to the forbidden
  set.
\item In the projection, $K\subseteq L$ makes the first term increase and
  the conditional $a$-cluster term decrease.
\item The disjoint transfer theorem is used only when
  $S\cap(Y\cup T)=\varnothing$.  If the source already contains a relay,
  Lemma~\ref{lem:source-overlap} supplies the missing branch.
\item In the edge step one must restrict first to $S\ncon a$.  Without
  that restriction, $C_a$ can change when $e$ is opened.
\item In Question~5, the separating vector is forced into the positive
  orthant because $r-\mathbb R_{\geq0}^{V}$ is unbounded in every
  negative coordinate direction.  This is what permits the use of an
  increasing linear cluster functional.
\item In Question~9, the random source is
  $S_\eta=\{o\}\cup N_\eta$, not merely $N_\eta$.  Retaining $o$ covers
  $o\in A$, $a=o$, and $b=o$ without exceptional cases.
\end{enumerate}

No limiting argument chooses a minimizer.  This is important at weights
$0$ or $1$ and in the presence of ties: the proof fixes the specified
minimizer $a$ throughout.

\section{Lean formalization plan}\label{sec:lean}

This section gives declarations at a granularity suitable for parallel
formalization.  Names are proposals; the mathematical interfaces are the
important part.

\subsection*{Block A: source clusters and event identities}

Reuse \lean{CovTau.sC} from
\lean{Percolation/Continuity/CovTau/A2Defs.lean}, or add a measure-level
wrapper \lean{openClusterSet} with mathematical specification
\[
 \operatorname{openClusterSet}(\omega,S)
 :=\{x:\exists s\in S,\ \operatorname{Reachable}_{\omega}(s,x)\}.
\]
Prove:
\begin{itemize}
\item singleton and union formulas;
\item monotonicity in the configuration and in $S$;
\item the pattern partition \eqref{eq:partition};
\item \lean{clusterSet\_pair\_independent\_internalEdge}, which is
  Lemma~\ref{lem:internal} in pointwise form.
\end{itemize}

\subsection*{Block B: avoided means and source surplus}

Define denominator-free numerators first:
\[
 M_x^Y=\EE[\Phi(C_x);x\ncon Y],\qquad d_x^Y=\PP(x\ncon Y).
\]
Define \lean{sourceFirstPattern} and \lean{sourceAvoidSurplus}.  Keep the
rank compatibility assumption in cross-multiplied form
\[
 M_x^Y d_z^Y\leq M_z^Y d_x^Y
\]
to avoid division at the boundary of the weight cube.

\subsection*{Block C: guarded CSH and corrected unfolding}

Formalize \eqref{eq:guardGamma}--\eqref{eq:guard-slack} first in
denominator-free form.  Make the whole-source exploration identity
\eqref{eq:guard-markov-a}, the correction \eqref{eq:guard-error}, and
the two facts \eqref{eq:guard-error-sign}--
\eqref{eq:guard-error-singleton} separate lemmas.  Prove the
cross-multiplied one-decoy identity \eqref{eq:guard-one-decoy-cross}
before the recursive unfolding \eqref{eq:guard-unfold}.  This order
prevents accidental reuse of the false identity with $C_d$ in place of
$C_R$.

Package the nonempty-intersection proof of \eqref{eq:guard-A2} as a
strong induction on the residual vertex set, exposing the cross-deleted
sets $P,P'$ in its four-functions input.  The two-half Gibbs step should
call the existing antitonicity and conditional-association declarations
before closing the regeneration contraction.  The output is
Theorem~\ref{thm:guardCSH}, first at nondegenerate weights and then in a
cross-multiplied boundary form.

\subsection*{Block D: set-source surplus transfer}

Generalize the outer observer in
\lean{CSH.surplusMargin\_nonneg\_of\_csh}.  The required repository pieces
already accept sets:
\begin{center}
\begin{tabular}{@{}ll@{}}
\toprule
role & existing declaration/file\\
\midrule
set cluster & \lean{CovTau.sC}, \lean{A2Defs.lean}\\
set marker covariance & \lean{CovTau.BfS}, \lean{StarBridgeS.lean}\\
two-source horizontal inequality & \lean{CovTau.a2H},
  \lean{CovTau.p1H\_univ}\\
two-set conditional association &
  \lean{BHK2006\_twoSetConditionalAssociation}\\
four-product exchange & \lean{setTwoClusterExchange}\\
ordinary CSH peeling & \lean{CSH/Peel.lean}\\
\bottomrule
\end{tabular}
\end{center}
The target declaration is the denominator-free form of
\eqref{eq:setS5}.  The proof follows the seven labelled stages in
\lean{CSH/Peel.lean}; insert $Y$ into every avoided set and use
$C_z\subseteq C_S$ at the source evaluation.

\subsection*{Block E: top-relay estimate and SAG}

Formalize Lemma~\ref{lem:toprelay} separately.  The Markov/tower step can
reuse the set-cluster exploration API; the monotonicity of $q(K)$ should
be its own lemma.  Then formalize the four displayed equations
\eqref{eq:surplus-peel}--\eqref{eq:kappa2} and close the induction with
Block~C.  Export both the ratio form \eqref{eq:SAG} (under positive
denominators) and a cross-multiplied boundary form.

\subsection*{Block F: set-source projection}

Generalize
\lean{CovTau.monotone\_projFun} and
\lean{CovTau.setIntegral\_sub\_eq\_projFun} from a singleton explored
cluster to \lean{openClusterSet}.  Prove precisely
\eqref{eq:projection-source} and \eqref{eq:projection-mean}.

\subsection*{Block G: common coefficients}

Formalize Theorem~\ref{thm:Q5} as a finite-dimensional alternative:
either a point of $\operatorname{conv}\{q^x:x\in A\}$ lies below $r$
coordinatewise, or a nonnegative separating functional contradicts
Conjecture~4 applied to $F_\lambda$.  A Mathlib theorem of alternatives
may replace strict separation, but the output must be one simplex vector
shared by every target $b$.

\subsection*{Block H: deleted-star conditioning}

Represent the incident-pair coordinates as a finite product and prove
\eqref{eq:star-source-identities} pointwise for each star state $\eta$.
Split on $S_\eta\con_H a$, prove the exact finite-sum identity
\eqref{eq:Q9-decomposition}, and apply the arbitrary-source version of
Theorem~\ref{thm:fixedmin}.  State the source as
$S_\eta=\{o\}\cup N_\eta$ to cover all overlap and endpoint cases.

\subsection*{Block I: final theorem statements}

Add named propositions for:
\begin{itemize}
\item \lean{kozmaNitzan\_conjecture4\_fixedMin};
\item \lean{kozmaNitzan\_conjecture2\_strong};
\item \lean{kozmaNitzan\_conjecture1};
\item the existing Conjecture~3 theorem as a corollary;
\item \lean{kozmaNitzan\_conjecture6\_strong}, whose statement omits the
  endpoint assumptions.
\item \lean{kozmaNitzan\_question5\_commonCoefficients};
\item \lean{kozmaNitzan\_question8};
\item \lean{kozmaNitzan\_question9}.
\end{itemize}

The recommended kernel audit is
\[
 \lean{\#print\ axioms kozmaNitzan\_conjecture6\_strong}
\]
followed by the corresponding commands for Conjectures 4, 2, and 1.
Question~7 is already represented by \lean{KNAll.question7\_holds};
Question~8 should be exported as a named one-line corollary of the
strong Conjecture~2 declaration.

\section{Conclusion}

The missing idea is not another estimate for a single observer.  It is to
retain the entire explored union $C_S$ when several possible observer
components are present.  The CSH proof was already built from set-cluster
exploration and two-set conditional association, so the source-set lift
preserves its induction.  With $S=\{o\}$ it selects the unconditional
minimizer in the pre-FKG comparison.  With $S=\{v,w\}$ it makes the state
of the edge $\{v,w\}$ internal to the source union, which is exactly the
invariance needed to force that edge open.  This proves all five numbered
Kozma--Nitzan conjectures.  Convex separation, the fixed-minimizer
specialization, and deleted-star conditioning then answer Questions
5 and 7--9.  Hence every numbered conjecture and question in the supplied
Kozma--Nitzan paper is resolved within the finite weighted-graph model.
The countable-graph closure then feeds Conjecture~3 into their reduction
and yields the paper's stated critical-percolation consequence.

\begin{thebibliography}{9}

\bibitem{KN}
G.~Kozma and S.~Nitzan,
\emph{A reduction of the $\theta(p_c)=0$ problem to a conjectured
inequality}, arXiv:2401.12397 (2024),
\url{https://arxiv.org/abs/2401.12397}.

\bibitem{BHK}
J.~van den Berg, O.~H\"aggstr\"om, and J.~Kahn,
\emph{Some conditional correlation inequalities for percolation and
related processes}, Random Structures Algorithms 29 (2006), 417--435.

\bibitem{Harris}
T.~E.~Harris,
\emph{A lower bound for the critical probability in a certain
percolation process}, Proc. Cambridge Philos. Soc. 56 (1960), 13--20.

\end{thebibliography}

\end{document}
